\section{Composition de l’holonomie \(G_9\) avec la carte \(T_1\) et certificat de précision arbitraire pour \(\alpha_{\rm HMT}\)}
\label{sec:alpha-t1-10000-rev9}

Le prolongement T1 applique la récurrence de retenue aux triades
\(P_k,E_k,\Phi_k\) produites préalablement par G9 et à la coordonnée
dodécaphasique déclarée \(K\) :
\begin{align}
 s_k&=P_k+E_k-\Phi_k-K_{k\bmod12}+c_{k+1},\\
 A_k&=s_k\bmod1000,\\
 c_k&=\left\lfloor\frac{s_k}{1000}\right\rfloor.
 \label{eq:composicion-g9-t1-alpha-10000-rev9}
\end{align}
Le domaine d’entrée contient les sorties ternaires engendrées par G9 ; il ne
contient ni un développement décimal de \(\alpha\), ni une valeur de
\(\alpha^{-1}\), ni CODATA, ni une condition métrologique cible.

Les cylindres de \(3501\) blocs fixent \(10020\) chiffres de garde dans chaque
canal, équivalant à \(3340\) triades. Le calcul propage, de droite à
gauche, toutes les conditions terminales possibles de l’ensemble stable
\[
 \mathcal C=\{-2,-1,0,1,2\}.
\]
Les cinq branches coalescent en un préfixe commun de \(3339\) triades, soit
\(10017\) chiffres. Par conséquent, les dix mille premiers chiffres de
\(\alpha_{\rm HMT}\) sont indépendants de la retenue située au-delà de la
zone de garde. La sortie publiée n’impose pas \(c_N=0\).

L’indépendance causale peut être vérifiée après la construction : les
trois mille premiers chiffres coïncident avec le développement décimal obtenu par le
lecteur externe, mais aucun d’eux n’intervient dans
\eqref{eq:composicion-g9-t1-alpha-10000-rev9}. Le développement de dix mille
chiffres et les données finies qui permettent de répéter le calcul figurent dans
l’annexe~\ref{app:desarrollos-10000-rev9}.
